% Generated draft: requires mathematical and source review.
\documentclass[11pt]{article}
\usepackage[margin=25mm]{geometry}
\usepackage{fontspec}
\setmainfont{Latin Modern Roman}
\usepackage{amsmath,amssymb,mathtools,unicode-math}
\setmathfont{Latin Modern Math}
\usepackage{enumitem,needspace,xurl,hyperref}
\setlength{\parindent}{0pt}
\setlength{\parskip}{5pt}
\setlength{\emergencystretch}{3em}
\newcommand{\N}{\mathbb{N}}
\newcommand{\Z}{\mathbb{Z}}
\newcommand{\Q}{\mathbb{Q}}
\newcommand{\R}{\mathbb{R}}
\newcommand{\C}{\mathbb{C}}
\newcommand{\F}{\mathbb{F}}
\newcommand{\A}{\mathbb{A}}
\newcommand{\PP}{\mathbb{P}}
\newcommand{\E}{\mathbb{E}}
\newcommand{\eps}{\varepsilon}
\newcommand{\dd}{\,\mathrm{d}}
\begin{document}
\section*{UnsolvedMath ID: 30006517}
\textbf{Source problem code:} catalogue entry
\subsection{Definitions and mathematical statement}
Equip the permutation group \(S_\infty=\operatorname{Sym}(\N)\) with the topology of pointwise convergence on the discrete set \(\N\). A permutation group \(G\leq S_\infty\) is \emph{oligomorphic} if, for every integer \(n\geq 1\), its coordinatewise action on \(\N^n\) has finitely many orbits.

For the classification question, adopt explicitly the usual closed permutation group interpretation: let \(\mathcal O\) be the space of closed oligomorphic subgroups of \(S_\infty\), with the Borel structure inherited from the Effros Borel space of closed subsets of \(S_\infty\). The Effros Borel structure is generated by the sets \(\{F:F\cap U\neq\varnothing\}\), where \(U\subseteq S_\infty\) is open. Closedness and this coding are conventions used to formulate the question here; they are not specified in the supplied excerpt and require confirmation against the intended source.

For \(G,H\in\mathcal O\), write \(G\cong_{\mathrm{top}}H\) if there is a group isomorphism \(\varphi:G\to H\) that is a homeomorphism. The topological isomorphism relation is \emph{smooth} if there exists a Borel map \(f:\mathcal O\to\R\) such that
\[
\forall G,H\in\mathcal O,\qquad
G\cong_{\mathrm{top}}H\quad\Longleftrightarrow\quad f(G)=f(H).
\]
\textbf{Classification question.} Does such a Borel map exist?

For a topological group \(G\), let \(\operatorname{Aut}_{\mathrm{top}}(G)\) be the group of group automorphisms of \(G\) that are homeomorphisms, and let
\[
\operatorname{Inn}(G)=\{x\mapsto gxg^{-1}:g\in G\},\qquad
\operatorname{Out}_{\mathrm{top}}(G)
=\operatorname{Aut}_{\mathrm{top}}(G)/\operatorname{Inn}(G).
\]
A topological group is \emph{profinite} if it is topologically isomorphic to an inverse limit of finite discrete groups. This property requires a specified topology, rather than only an abstract group structure.

\textbf{Related existence question.} Does there exist an oligomorphic group \(G\) whose outer automorphism group, with the topology intended by the source, is not profinite? Under the closed permutation group interpretation above, the domain is \(G\in\mathcal O\). The supplied record does not specify whether its automorphisms are required to be continuous, nor does it specify a topology on the automorphism group or on its outer quotient. Consequently, the displayed definition using topological automorphisms is an explicit interpretation, and the profiniteness question remains conditional on clarification of the intended automorphism and topology conventions.

The supplied passage attributes to Nies, Schlicht, and Tent the result that topological isomorphism of oligomorphic groups is Borel reducible to a Borel equivalence relation with countable classes. Here a Borel reduction from \(E\) on \(X\) to \(F\) on \(Y\) means a Borel map \(r:X\to Y\) satisfying \(x\mathrel E x'\) if and only if \(r(x)\mathrel F r(x')\). The passage presents the outer automorphism question as important to the smoothness question; it supplies no precise implication or equivalence between their answers.
\subsection{Short English statement}
Can closed oligomorphic permutation groups be assigned Borel real invariants that agree exactly when the groups are topologically isomorphic? The source also asks whether an oligomorphic group can have a nonprofinite outer automorphism group. That second question requires clarification of the intended automorphisms and topology. Closedness and the Borel coding are explicit conventions in this draft. The supplied dated assessment labels the problem partially solved, with the general questions remaining open; this assessment is not independently verified.
\subsection{Sources}
\begin{itemize}
\item Local UnsolvedMath registry, supplied complete record: permanent ID \(30006517\), exact source problem code \texttt{catalogue entry}, title \emph{Smooth Classification and Outer Automorphisms of Oligomorphic Groups}; collection: Oberwolfach Reports Open Problems. No individual proposer is specified; the recorded proposed year is \(2025\).
\item Supplied report citation: \emph{Homogeneous Structures: Model Theory meets Universal Algebra (2026)}, Oberwolfach Reports, pp.~3063--3124, DOI \texttt{10.4171/owr/2025/58}. \url{https://doi.org/10.4171/owr/2025/58}. The registry describes the original passage as retained verbatim from workshop open-problem discussions. The supplied citation uses the year \(2026\), whereas the DOI contains \(2025\); this discrepancy is retained without resolution.
\item Paolini--Nies, \emph{Two classes of oligomorphic groups with smooth isomorphism relation}. \url{https://arxiv.org/abs/2410.02248}.
\item Nies--Schlicht--Tent, \emph{Coarse groups and oligomorphic groups}. \url{https://arxiv.org/abs/1903.08436}. The supplied passage attributes to these authors the reduction to a Borel equivalence relation with countable classes.
\item Supplied literature assessment, checked on exactly \texttt{2026-08-22}: status \texttt{partially\_solved}. It reports that topological-isomorphism smoothness is proved for important subclasses, that their outer automorphism groups are profinite, and that the general smoothness question and existence of a nonprofinite example remain open. This is supplied dated curation derived from the OWR record, not an independently verified literature assessment; the subclasses and their hypotheses are not identified in the supplied record.
\item The registry labels the statement \texttt{corrected\_verified} and reports that the official report was checked, the question isolated from surrounding exposition, and notation repaired for self-contained publication. These are supplied verification claims, not verification performed for this draft.
\end{itemize}
\end{document}
